% !PRINTED 101
\weberpage{101}

\begin{center}
{\Large\gesperrt{Dritter Abschnitt.}}

\vspace{1em}
{\large\gesperrt{Die Wurzeln algebraischer Gleichungen.}}
\end{center}

\begin{center}
{\large \S.\ 30.}

\vspace{0.5em}
{\large\gesperrt{Begriff der Wurzeln. Mehrfache Wurzeln.}}
\end{center}

Nachdem in den beiden ersten Abschnitten die algebraischen Grössen mehr von der formalen Seite betrachtet waren, wobei es sich um identische Umformungen von Buchstabenausdrücken handelte, in denen die Buchstaben durchweg als Symbole für variable Grössen aufgefasst werden konnten, treten nun die Zahlengrössen mehr in den Vordergrund.

Wir verstehen hier unter Zahlen, gemäss dem in der Einleitung Festgesetzten, reelle oder imaginäre Grössen von der Form $a+bi$, und stellen die reellen Grössen zur Veranschaulichung durch die Punkte einer geraden Linie, die imaginären durch die Punkte einer Ebene dar. Unter dem \gesperrt{absoluten Werth} einer imaginären Grösse $a+bi$ verstehen wir $\sqrt{a^2+b^2}$ und bezeichnen ihn nach Weierstrass mit
\[
 |a+bi|.
\]
Es sei nun
\begin{equation}
 f(x)=a_0x^n+a_1x^{n-1}+a_2x^{n-2}+\cdots+a_n
\tag{1}
\end{equation}
eine ganze Function von $x$, worin die Coëfficienten irgend welche reelle oder imaginäre Zahlen sind, und der erste, $a_0$, von Null verschieden vorausgesetzt wird. Wenn $\alpha$ eine Zahl ist, die für $x$ gesetzt die Function $f(x)$ zu Null macht, die also der Bedingung $f(\alpha)=0$ genügt, so heisst $\alpha$ eine Wurzel der Gleichung
\[
 f(x)=0.
\]
